\chapter{Computational Results}
\label{chap:results}

\vspace*{\fill}

This chapter evaluates the hot-started nullspace active-set method on generated instances of
the \textbf{airline crew scheduling problem}. Every instance is solved in two variants: with
equality coverage rows only, and with 30\% of the coverage rows relaxed to ``$\ge$''
inequalities.

\medskip

We first check \textbf{correctness}: every column generation method must reproduce the
optimum of the full linear program, which is still small enough to be solved directly. We
then study the \textbf{convergence behaviour} of regularized column generation, which
consists of two clearly separated kinds of iterations, and the \textbf{computational time},
broken down by solver, by iteration and by operation inside the active-set method.

\medskip

The last two sections look inside the method. We follow the sizes of the nullspace sets
over a run with inequality rows, and we measure the \textbf{contents of the factorization}:
how many of the stored entries are actually nonzero, how fast structural zeros accumulate,
and how much periodic refactorization is worth.

\vspace*{\fill}

\newpage


\section{Test Problem: Airline Crew Scheduling}
\label{sec:test-problem}

\subsection{The model}

In the crew scheduling problem, $K$ pilots have to be assigned schedules that together
operate a set of $m$ flights (\cref{ex:crew}). The schedules of a pilot are the source--sink
paths in a \textbf{time--space network}: a directed acyclic graph $G = (V,E)$ whose nodes
represent states of a pilot at given times, and whose edges represent the allowed transitions
between them. Edge $e$ costs pilot $k$ the amount $c_{e,k}$, and it operates a subset of the
flights, recorded in the \textbf{task matrix} $A_{\mathrm{task}}\in\{0,1\}^{m\times|E|}$. With
$x_k\in\{0,1\}^{|E|}$ the edge vector of the path of pilot $k$ and $F$ the node--edge incidence
matrix of $G$, the linear relaxation reads
\begin{subequations} \label{eq:crew}
\begin{align}
\min_{x_1,\dots,x_K}\quad & \sum_{k=1}^{K} c_k^\top x_k\\
\text{s.t.}\quad & \Bigl(\sum_{k=1}^{K}A_{\mathrm{task}}\,x_k\Bigr)_i = 1 && (i\in\mathcal{E}),
\label{eq:crew-eq}\\
& \Bigl(\sum_{k=1}^{K}A_{\mathrm{task}}\,x_k\Bigr)_i \ge 1 && (i\in\mathcal{I}),
\label{eq:crew-ineq}\\
& F x_k = d,\quad 0\le x_k\le 1 && (k = 1,\dots,K),
\end{align}
\end{subequations}
where $d$ has the entry $-1$ at the source, $+1$ at the sink and $0$ elsewhere. This is a
multi-block linear program (\cref{def:multiblock}) with identical block polytopes. A
coverage row in $\mathcal{E}$ requires a flight to be operated exactly once; a row in
$\mathcal{I}$ requires it at least once, which allows a pilot to travel on a flight as a
passenger. The columns of the master problem are complete schedules, and the pricing
problem~\eqref{eq:pricing} of pilot $k$ is a shortest path problem with edge weights
$c_{e,k} - (A_{\mathrm{task}}^\top\pi)_e$. These weights can be negative, but the graph is
acyclic, so one pass of dynamic programming in topological order solves it in
$O(|V|+|E|)$.

\subsection{Instance generator}

The instances are generated randomly but reproducibly.
\begin{itemize}
    \item \textbf{Network.} The nodes $1,\dots,|V|$ are ordered in time. A backbone edge
          $v\to v+1$ connects consecutive nodes, which guarantees a source--sink path. Every
          skip edge $v\to w$ with $w\ge v+2$ exists independently with a fixed probability. An
          edge costs $(w-v) + U(0,1)$, multiplied by $1 + 0.1\,U(0,1)$ for each pilot
          separately.
    \item \textbf{Flights.} Each skip edge operates a random set of up to ten flights whose
          position in time is close to that of the edge. A flight is only attached to edges
          that overlap in time, so that a single schedule operates every flight at most once.
          The task matrix therefore has a banded structure.
    \item \textbf{Feasibility.} Only flights that can be covered exactly once in the linear
          relaxation are kept, so that the equality variant is feasible.
    \item \textbf{Inequality rows.} In the inequality variant, a random 30\% of the coverage
          rows are turned into ``$\ge1$'' rows.
\end{itemize}

\Cref{tab:instances} lists the four instance sizes. The last column is the number of
variables of the full edge-based formulation~\eqref{eq:crew}, which HiGHS solves directly to
provide the reference optimum $z^\star$.

\begin{table}[ht]
\centering
\renewcommand{\arraystretch}{1.2}
\begin{tabular}{lrrrrrr}
\toprule
instance & $|V|$ & $|E|$ & $K$ & $m$ & $m_I$ (inequality variant) & variables of~\eqref{eq:crew} \\
\midrule
\tabinput{data/tab_instances.tex}
\bottomrule
\end{tabular}
\caption{Test instances. $K$ is the number of pilots, $m$ the number of coverage rows, and
$m_I$ the number of ``$\ge$'' rows in the inequality variant.}
\label{tab:instances}
\end{table}


\section{Implementation and Setup}
\label{sec:setup}

All methods are implemented in Julia~1.12 \cite{bezanson2017julia} in one framework: the
column generation driver, the pricing algorithm, the column pool and the statistics are
shared, and only the solver of the restricted master problem differs. We compare five
methods:

\begin{center}
\renewcommand{\arraystretch}{1.2}
\resizebox{\textwidth}{!}{%
\begin{tabular}{ll}
\toprule
method & master problem \\
\midrule
Exact LP (HiGHS) & none: formulation~\eqref{eq:crew} solved directly by the HiGHS simplex method \\
LP CG (HiGHS) & classical linear master problem, HiGHS simplex, warm-started \\
Explicit QP CG (HiGHS) & explicit regularized master problem~\eqref{eq:explicit}, HiGHS QP solver \\
Eliminated QP CG (HiGHS) & eliminated regularized master problem~\eqref{eq:eliminated}, HiGHS QP \\
Hot-started QPAS CG & eliminated regularized master problem, \cref{alg:hot-qpas} \\
\bottomrule
\end{tabular}
}
\end{center}

HiGHS \cite{huangfu2018parallelizing} (version 1.15) is used through a persistent model,
to which new columns are added incrementally, so that its LP solves are warm-started. The
eliminated QP is only defined for equality rows in HiGHS form, since it requires the
elimination of $w$ without surplus variables; it is therefore run on the equality variant
only. The active-set method uses CHOLMOD 5.3 \cite{chen2008cholmod} for the factorization and
its updates.

\medskip

All column generation runs use the same settings: the initial columns are the cheapest path
of every pilot, pricing visits the pilots in random order and adds the first improving
column, the proximal point rule moves the center after every iteration, and $\delta = 1$,
$\tau = 10^{-6}$, $\varepsilon_{\mathrm{rc}} = 10^{-6}$. The active-set method uses the
column capacity $\nmax = 20\,000$ and the refactorization interval $150$ unless stated
otherwise. The HiGHS QP masters are stopped after 600 seconds and the LP master after 900
seconds. All experiments ran single-threaded on an AMD Ryzen~7 7700 processor with 32\,GB of
memory under Windows~11. Times of the active-set method are the faster of two runs.


\section{Correctness}
\label{sec:results-correctness}

\Cref{tab:comparison-eq,tab:comparison-ineq} compare all methods on both variants. Besides the
status, the number of column generation iterations, the number of generated columns and the
times, they report three accuracy measures relative to the reference optimum $z^\star$:
\begin{itemize}
    \item the relative error of the cost, $(c^\top\lambda - z^\star)/|z^\star|$;
    \item the relative gap of the best Lagrangian bound, $(z^\star - \theta)/|z^\star|$, which
          is nonnegative for every valid bound (\cref{prop:lagrangian-bound});
    \item the infeasibility $\lVert b - A\lambda\rVert_\infty$ of the final weights in the linking
          rows (for inequality rows only the violated part).
\end{itemize}

\begin{table}[p]
\centering
\scriptsize
\renewcommand{\arraystretch}{1.15}
\setlength{\tabcolsep}{4pt}
\resizebox{\textwidth}{!}{%
\begin{tabular}{llcrrrrccc}
\toprule
 & method & status & CG iter. & columns & time [s] & master [ms/iter.] & cost error & bound gap & infeasibility \\
\midrule
\tabinput{data/tab_comparison_eq.tex}
\bottomrule
\end{tabular}
}
\caption{Equality variant: all methods. ``time'' denotes a method stopped by its time limit.
Master times are averages per column generation iteration.}
\label{tab:comparison-eq}

\bigskip\bigskip

\resizebox{\textwidth}{!}{%
\begin{tabular}{llcrrrrccc}
\toprule
 & method & status & CG iter. & columns & time [s] & master [ms/iter.] & cost error & bound gap & infeasibility \\
\midrule
\tabinput{data/tab_comparison_ineq.tex}
\bottomrule
\end{tabular}
}
\caption{Inequality variant (30\% ``$\ge$'' rows). The eliminated HiGHS QP master does not
support inequality rows.}
\label{tab:comparison-ineq}
\end{table}

\paragraph{All methods reach the optimum.} Every column generation run that finishes within
its time limit stops with an optimal status, and the Lagrangian bound agrees with $z^\star$ to
within $10^{-6}$ relative or better. This confirms the exact termination criterion of
\cref{thm:exact-termination} in practice, including the sign conditions of the inequality
rows: the bound would be $-\infty$ if the projected duals violated them
(\cref{rem:bound-sign}).

\paragraph{The accuracy of the stabilized methods.} The stabilized methods stop when
$\lVert\pi-\pihat\rVert_\infty\le\tau\max(1,\lVert\pi\rVert_\infty)$, which by
\cref{prop:stop-accuracy} allows an infeasibility of up to $\delta\tau\max(1,\lVert\pi\rVert_\infty)$
in the linking rows. The infeasibility column shows exactly this order of magnitude, about
$10^{-5}$ for the larger instances, where the duals are of order $10$ to $100$. The cost
$c^\top\lambda$ of such a slightly infeasible solution can be off by more than the
infeasibility itself: on I2 it is $2\cdot10^{-5}$ below $z^\star$. The Lagrangian bound, which
certifies the optimal value, is unaffected and agrees with $z^\star$ to about $10^{-7}$. A
tighter $\tau$ removes the cost error at the price of a few more center-only iterations. In a
test on a smaller instance with $\delta = 5$, the cost error dropped from $10^{-6}$ for
$\tau = 10^{-6}$ to $10^{-15}$ for $\tau = 10^{-9}$, with one additional iteration.

\paragraph{The active-set method and HiGHS solve the same problem.} On instance I1, where
all stabilized methods finish, the explicit and eliminated HiGHS QPs and the active-set method
perform almost the same number of column generation iterations and generate the same number
of columns: they follow the same sequence of master problems, as they should. The small
differences come from the accuracy of the master solves, which affects the duals passed to
pricing and hence the order in which columns are found.


\section[Convergence and the Two Phases]{Convergence and the Two Phases of Regularized Column Generation}
\label{sec:results-phases}

\Cref{fig:convergence} shows the convergence of the master objective on instance I2 for all
column generation methods, together with the change of the duals from one iteration to the
next.

\begin{figure}[ht]
\centering
\begin{tikzpicture}
\begin{groupplot}[group style={group size=2 by 1, horizontal sep=1.6cm},
    width=0.5\textwidth, height=0.42\textwidth, ymode=log,
    xlabel={CG iteration}, tick label style={font=\scriptsize}, label style={font=\small},
    title style={font=\small}, legend style={font=\scriptsize, at={(0.98,0.98)}, anchor=north east},
    every axis plot/.append style={line width=0.7pt}]
\nextgroupplot[title={relative error of the master objective}, ymin=1e-9, ymax=1]
    \addplot[cred] table[x=it, y=gap, col sep=comma] {data/conv_I2_0_lp.csv};
    \addplot[cgreen] table[x=it, y=gap, col sep=comma] {data/conv_I2_0_explicit_qp.csv};
    \addplot[cpurple, densely dashed] table[x=it, y=gap, col sep=comma] {data/conv_I2_0_eliminated_qp.csv};
    \addplot[cblue] table[x=it, y=gap, col sep=comma] {data/conv_I2_0_qpas.csv};
    \legend{LP CG, explicit QP, eliminated QP, QPAS}
\nextgroupplot[title={dual change $\lVert\pi_t - \pi_{t-1}\rVert_\infty$}, ymin=1e-6]
    \addplot[cred, only marks, mark size=0.3pt] table[x=it, y=dpi, col sep=comma] {data/conv_I2_0_lp.csv};
    \addplot[cblue, only marks, mark size=0.3pt] table[x=it, y=dpi, col sep=comma] {data/conv_I2_0_qpas.csv};
    \legend{LP CG, QPAS}
\end{groupplot}
\end{tikzpicture}
\caption{Instance I2 (equality variant). Left: relative distance of the master objective to
$z^\star$. Right: change of the duals between consecutive iterations. The two HiGHS QP masters
stop at their time limit after about 6\,000 iterations.}
\label{fig:convergence}
\end{figure}

\paragraph{Stabilization.} The duals of the linear master problem jump by up to $2\cdot10^4$
in the first iterations, when the artificial Big-M variables are still active, and by about
$2$ in a typical iteration of the second half of the run. The regularized duals change by
about $10^{-4}$ in a typical iteration of the second half. The linear master problem, however,
reaches the optimum in fewer iterations on this instance. Its jumps explore the dual space
quickly, while the proximal term keeps the regularized duals close to the center, so that
they need many small steps. Regularization therefore does not reduce the number of
iterations here. What it buys is iterations whose master problems differ only slightly from
each other, which is exactly what the hot start exploits: an iteration of the active-set
method costs $0.2$\,ms on average, against $30$\,ms for the warm-started LP master
(\cref{tab:comparison-eq}).

\paragraph{The QP masters follow the same path.} The curves of the explicit QP, the
eliminated QP and the active-set method are almost indistinguishable over the first 6\,000
iterations, after which the HiGHS QP masters run out of time.

\medskip

The right panel of \cref{fig:convergence} shows a second feature of the regularized methods:
from some point on, the duals change by very small amounts in most iterations. These are the
center-only iterations of \cref{sec:center-update}. \Cref{fig:phases} shows how the two kinds
of iterations are distributed over a run.

\begin{figure}[ht]
\centering
\begin{tikzpicture}
\begin{groupplot}[group style={group size=2 by 2, horizontal sep=1.6cm, vertical sep=0.6cm,
    xlabels at=edge bottom, xticklabels at=edge bottom},
    width=0.5\textwidth, height=0.3\textwidth, xlabel={CG iteration},
    tick label style={font=\scriptsize}, label style={font=\small}, title style={font=\small},
    every axis plot/.append style={line width=0.7pt}]
\nextgroupplot[title={I2, equality variant}, ylabel={columns}]
    \addplot[cblue] table[x=it, y=ncol, col sep=comma] {data/conv_I2_0_qpas.csv};
\nextgroupplot[title={I2, inequality variant}]
    \addplot[cblue] table[x=it, y=ncol, col sep=comma] {data/conv_I2_30_qpas.csv};
\nextgroupplot[ylabel={share adding a column}, ymin=0, ymax=1.05]
    \addplot[const plot mark mid, fill=cblue!35, draw=cblue] table[x=it, y=share, col sep=comma]
        {data/phases_I2_0.csv} \closedcycle;
\nextgroupplot[ymin=0, ymax=1.05]
    \addplot[const plot mark mid, fill=cblue!35, draw=cblue] table[x=it, y=share, col sep=comma]
        {data/phases_I2_30.csv} \closedcycle;
\end{groupplot}
\end{tikzpicture}
\caption{The two phases of regularized column generation (active-set method, instance I2).
Top: number of columns. Bottom: share of the iterations in each window of 100 iterations
that add a column; the remaining iterations are center-only iterations.}
\label{fig:phases}
\end{figure}

\paragraph{Two phases.} A run starts with a \textbf{column phase}, in which (almost) every
iteration adds a column. It is followed by a \textbf{proximal phase}, in which most
iterations are center-only iterations that move the center towards the restricted dual
optimum, interrupted by occasional iterations that find a new column because the duals have
moved. The phases are not strictly separated: columns keep appearing until close to the end.
\Cref{tab:phases} quantifies this for all instances.

\begin{table}[ht]
\centering
\small
\renewcommand{\arraystretch}{1.15}
\setlength{\tabcolsep}{4pt}
\resizebox{\textwidth}{!}{%
\begin{tabular}{lcrrrrrrrr}
\toprule
 & & & & & first & last & \multicolumn{2}{c}{mean inner iter.} & master time \\
instance & ineq. & CG iter. & with column & center-only & center-only & column & column & center-only & center-only \\
\midrule
\tabinput{data/tab_phases.tex}
\bottomrule
\end{tabular}
}
\caption{Column-adding and center-only iterations of the active-set method. ``inner iter.''
counts the iterations of the active-set method per master problem.}
\label{tab:phases}
\end{table}

The two kinds of iterations differ sharply in cost. A column-adding iteration needs on
average tens of active-set iterations, since the new column is released and the working set
adapts to it. A center-only iteration needs at most two on average: as shown in
\cref{prop:center-only}, once the working set has settled, the direction is zero, and the
master problem is solved by a single multiplier test. The length of the proximal phase
varies a lot between instances. On I2 and I3 (equality variant) center-only iterations are
the large majority, on I4 they are a minority. In all cases their share of the master time
is much smaller than their share of the iterations.

\Cref{fig:optimality} shows how the optimality measures evolve over the two phases for the
active-set method on I2. In the column phase, the cost error, the error of the master
objective and the infeasibility decrease together. In the proximal phase, the Lagrangian
bound becomes available, since all pilots are priced in every center-only iteration. The
bound approaches $z^\star$ from below at the same rate at which the master objective
approaches it, so the two errors are indistinguishable in the plot. The infeasibility, and
with it the cost error, decrease more slowly and in steps: each step is a change of the
working set that ends one of the dual walks of \cref{prop:center-only}.

\begin{figure}[ht]
\centering
\begin{tikzpicture}
\begin{axis}[width=0.78\textwidth, height=0.42\textwidth, ymode=log, xlabel={CG iteration},
    ymin=1e-9, ymax=10, tick label style={font=\scriptsize}, label style={font=\small},
    legend style={font=\scriptsize, at={(0.98,0.98)}, anchor=north east, legend columns=2},
    every axis plot/.append style={line width=0.7pt}]
    \addplot[cred!45, only marks, mark=*, mark size=0.7pt] table[x=it, y=bgap, col sep=comma] {data/bound_I2_0_qpas.csv};
    \addplot[corange] table[x=it, y=costgap, col sep=comma] {data/conv_I2_0_qpas.csv};
    \addplot[cgreen] table[x=it, y=infeas, col sep=comma] {data/conv_I2_0_qpas.csv};
    \addplot[cblue, densely dashed] table[x=it, y=gap, col sep=comma] {data/conv_I2_0_qpas.csv};
    \legend{bound gap $(z^\star-\theta)/|z^\star|$, cost error $|c^\top\lambda - z^\star|/|z^\star|$,
            infeasibility, master objective error}
\end{axis}
\end{tikzpicture}
\caption{Optimality measures of the active-set method over a run (I2, equality variant). The
Lagrangian bound is only available in iterations that price all pilots, which are the
iterations that find no column.}
\label{fig:optimality}
\end{figure}


\section{Computational Time}
\label{sec:results-time}

\paragraph{Master time per iteration.} \Cref{fig:master-time} shows the time of the master
solve per column generation iteration on I2. The time of the LP master and of the two HiGHS QP
masters grows with the number of columns. The HiGHS QP solvers take 100--170\,ms per
iteration in the second half of the run, since they rebuild their internal factorization for
every master problem. The active-set method stays between $0.01$ and $2$\,ms per iteration
over the whole run, with a median of $0.08$\,ms in the second half. That is about three orders
of magnitude below the HiGHS QP masters and more than two orders of magnitude below the
warm-started LP master.

\begin{figure}[ht]
\centering
\begin{tikzpicture}
\begin{axis}[width=0.78\textwidth, height=0.42\textwidth, ymode=log, xlabel={CG iteration},
    ylabel={master time per iteration [ms]}, tick label style={font=\scriptsize},
    label style={font=\small}, legend style={font=\scriptsize, at={(0.98,0.5)}, anchor=east},
    every axis plot/.append style={line width=0.8pt}]
    \addplot[cred] table[x=it, y=mtavg, col sep=comma] {data/conv_I2_0_lp.csv};
    \addplot[cgreen] table[x=it, y=mtavg, col sep=comma] {data/conv_I2_0_explicit_qp.csv};
    \addplot[cpurple, densely dashed] table[x=it, y=mtavg, col sep=comma] {data/conv_I2_0_eliminated_qp.csv};
    \addplot[cblue] table[x=it, y=mtavg, col sep=comma] {data/conv_I2_0_qpas.csv};
    \legend{LP CG, explicit QP, eliminated QP, QPAS}
\end{axis}
\end{tikzpicture}
\caption{Master time per column generation iteration on I2 (moving average over 50
iterations).}
\label{fig:master-time}
\end{figure}

\paragraph{Where the time goes.} \Cref{fig:breakdown} splits the time of the active-set
method into its operations. The largest share, 37--60\%, is the direction computation (the
right-hand side and the sparse solve with the factor). On the larger instances it is followed
by the updates of the factorization and the periodic refactorization, together 35--50\%. The
multiplier test, the only operation whose cost grows with the number of columns, takes
10--17\%. Pricing takes only 1--8\%, except on the tiny instance I1.

\begin{figure}[ht]
\centering
\begin{tikzpicture}
\begin{axis}[xbar stacked, table/col sep=comma, width=0.8\textwidth, height=0.5\textwidth, xmin=0, xmax=100,
    xlabel={share of the total time [\%]}, ytick=data,
    yticklabels from table={data/breakdown.csv}{label},
    y dir=reverse, bar width=9pt, tick label style={font=\scriptsize}, label style={font=\small},
    legend style={font=\scriptsize, at={(0.5,-0.18)}, anchor=north, legend columns=3},
    reverse legend=false]
    \addplot[fill=cblue!70, draw=cblue] table[x=direction, y expr=\coordindex, col sep=comma] {data/breakdown.csv};
    \addplot[fill=corange!70, draw=corange] table[x=multipliers, y expr=\coordindex, col sep=comma] {data/breakdown.csv};
    \addplot[fill=cgrey!50, draw=cgrey] table[x=step, y expr=\coordindex, col sep=comma] {data/breakdown.csv};
    \addplot[fill=cgreen!60, draw=cgreen] table[x=updates, y expr=\coordindex, col sep=comma] {data/breakdown.csv};
    \addplot[fill=cred!60, draw=cred] table[x=refactor, y expr=\coordindex, col sep=comma] {data/breakdown.csv};
    \addplot[fill=cpurple!50, draw=cpurple] table[x=pricing, y expr=\coordindex, col sep=comma] {data/breakdown.csv};
    \legend{direction, multipliers, step, updates, refactorization, pricing}
\end{axis}
\end{tikzpicture}
\caption{Time breakdown of the active-set method and of pricing, as shares of their sum.
``updates'' are the expansions and reductions of the working set without the
refactorizations.}
\label{fig:breakdown}
\end{figure}

\Cref{tab:counters} lists the events inside the active-set method for every run. Over a full
run, the working set changes tens of thousands of times, and every change is a single update
of the factorization. Representative swaps are rare, because representatives carry the
largest weights. Singular reduced Hessians occur, but only a handful of times per run, and
each is resolved by one zero-curvature step. The final nullspace dimension $|\Nx|$ stays below
the number of coverage rows, as guaranteed by \cref{cor:dimension}.

\begin{table}[ht]
\centering
\scriptsize
\renewcommand{\arraystretch}{1.15}
\setlength{\tabcolsep}{3.5pt}
\resizebox{\textwidth}{!}{%
\begin{tabular}{lcrrrrrcrrrrr}
\toprule
 & & AS iter. & steps (full) & expans. & swaps & reduct. & surplus exp./red. & sing. & updates & refact. & $|\Nx|$ & $|\Ns|$ \\
\midrule
\tabinput{data/tab_qpas_counters.tex}
\bottomrule
\end{tabular}
}
\caption{Events inside the active-set method over a whole column generation run: active-set
iterations, primal steps (of which full steps), expansions, representative swaps, reductions,
expansions/reductions of surpluses, singular steps, factorization updates, refactorizations, and
the final sizes of $\Nx$ and $\Ns$.}
\label{tab:counters}
\end{table}

\paragraph{Scaling.} \Cref{fig:scaling} shows the total time of the active-set method and of
the direct LP solve as the instances grow, up to an instance with 960 coverage rows
(instances S5 and S6 with $|V| = 300$ and $400$, $K = 16$ and $20$, generated in the same
way). Over the whole range, both grow roughly with the third power of $m$. For the active-set
method this is expected: $|\Nx|$ grows proportionally to $m$, the factor of $\Heff$ is dense
(\cref{sec:results-fill}), and the number of iterations grows as well. From I4 on, the
complete column generation run with the active-set method is faster than the direct solution
of the full LP. On S6 the direct LP solve did not finish within 20 minutes, while the
active-set method solved both variants in 4--7 minutes. The inequality variant is
consistently faster than the equality variant, since it needs fewer iterations
(\cref{sec:results-ineq}).

\begin{figure}[ht]
\centering
\begin{tikzpicture}
\begin{axis}[width=0.7\textwidth, height=0.48\textwidth, xmode=log, ymode=log,
    xlabel={coverage rows $m$}, ylabel={total time [s]}, tick label style={font=\scriptsize},
    label style={font=\small}, legend style={font=\scriptsize, at={(0.02,0.98)}, anchor=north west},
    every axis plot/.append style={line width=0.8pt}, log basis x=10]
    \addplot[cblue, mark=*, mark size=1.5pt] table[x=m, y=qpas, col sep=comma] {data/scaling_eq.csv};
    \addplot[cblue, dashed, mark=o, mark size=1.5pt] table[x=m, y=qpas, col sep=comma] {data/scaling_ineq.csv};
    \addplot[cred, mark=square*, mark size=1.5pt] table[x=m, y=exact, col sep=comma] {data/scaling_eq.csv};
    \addplot[cred, dashed, mark=square, mark size=1.5pt] table[x=m, y=exact, col sep=comma] {data/scaling_ineq.csv};
    \addplot[cgrey, dotted, domain=40:1000, samples=2] {4e-7*x^3};
    \legend{QPAS CG (equality), QPAS CG (inequality), exact LP (equality), exact LP (inequality), $\propto m^3$}
\end{axis}
\end{tikzpicture}
\caption{Total time against the number of coverage rows for instances I1--I4, S5 and S6. The
exact LP of S6 did not finish within 20 minutes.}
\label{fig:scaling}
\end{figure}


\section{Inequality Rows}
\label{sec:results-ineq}

With 30\% inequality rows, column generation needs considerably fewer iterations on all
instances (\cref{tab:comparison-eq,tab:comparison-ineq}): over-covering a flight is allowed,
so the optimal schedules are found with fewer columns, and the proximal phase is shorter.

\medskip

\Cref{fig:sets} follows the sizes of the nullspace sets over a run with inequality rows.
Within the first iterations, about 35 surpluses are released: the initial schedules
over-cover their rows, so the duals of these rows become negative. These rows drop out of
$\Aeff$, and by \cref{cor:dimension} they also reduce the room for nullspace columns,
$|\Nx|\le m - |\Ns|$. Over the rest of the run $|\Ns|$ only changes slowly, reaching 48, while
individual surpluses are released and fixed again many times (\cref{tab:counters}). The number
of nullspace columns grows during the column phase and ends at $|\Nx| = 306$, just below its
bound $m - |\Ns| = 308$: at the end of the run, the effective reduced Hessian is almost as
large as it can be.

\begin{figure}[ht]
\centering
\begin{tikzpicture}
\begin{axis}[width=0.78\textwidth, height=0.42\textwidth, xlabel={CG iteration},
    ylabel={set size}, tick label style={font=\scriptsize}, label style={font=\small},
    legend style={font=\scriptsize, at={(0.98,0.5)}, anchor=east},
    every axis plot/.append style={line width=0.8pt}, ymin=0]
    \addplot[cgrey, dashed] table[x=it, y=m, col sep=comma] {data/sets_I3_30.csv};
    \addplot[cred] table[x=it, y=bound, col sep=comma] {data/sets_I3_30.csv};
    \addplot[cblue] table[x=it, y=N, col sep=comma] {data/sets_I3_30.csv};
    \addplot[corange] table[x=it, y=Ns, col sep=comma] {data/sets_I3_30.csv};
    \legend{$m$, $m - |\Ns|$ (bound on $|\Nx|$), $|\Nx|$, $|\Ns|$}
\end{axis}
\end{tikzpicture}
\caption{Nullspace columns and free surpluses over a run (I3, inequality variant), sampled
every 20 iterations.}
\label{fig:sets}
\end{figure}

\medskip

Surplus operations are frequent (\cref{tab:counters}): with inequality rows, 13--19\% of all
expansions and reductions of the working set on I2--I4 are expansions and reductions of
surpluses, each of which is a rank-one update or downdate of the factorization. They do not
make the method slower: the average time per master problem of the inequality variant is of
the same order as that of the equality variant (\cref{tab:comparison-eq,tab:comparison-ineq}).
The remaining differences come from the mix of iterations. On I2, for instance, the inequality
variant has relatively few of the cheap center-only iterations (\cref{tab:phases}), which raises
its average. Singular
reduced Hessians, of columns or surpluses, occurred at most eight times per run. The
singular surplus release of \cref{thm:zc-surplus} is exercised by the test suite of the
implementation on constructed instances.


\section{Fill-In and Refactorization}
\label{sec:results-fill}

The factorization of $\Heff$ is updated tens of thousands of times per run
(\cref{tab:counters}). \Cref{fig:fill} shows what this does to its contents, using the
quantities of \cref{def:stored-effective}, without refactorization (left) and with
refactorization every 150 updates (right).

\begin{figure}[ht]
\centering
\begin{tikzpicture}
\begin{groupplot}[group style={group size=2 by 1, horizontal sep=1.6cm},
    width=0.5\textwidth, height=0.42\textwidth, xlabel={CG iteration}, ymode=log,
    ymin=100, ymax=2e6, tick label style={font=\scriptsize}, label style={font=\small},
    title style={font=\small}, legend style={font=\scriptsize, at={(0.98,0.02)}, anchor=south east},
    every axis plot/.append style={line width=0.7pt}]
\nextgroupplot[title={never refactorized}, ylabel={off-diagonal entries of $L$}]
    \addplot[cred] table[x=it, y=stored, col sep=comma] {data/fill_I3_tau0.csv};
    \addplot[cblue] table[x=it, y=effective, col sep=comma] {data/fill_I3_tau0.csv};
    \addplot[cgrey, dashed] table[x=it, y=dense, col sep=comma] {data/fill_I3_tau0.csv};
    \legend{stored, effective (nonzero), dense equivalent}
\nextgroupplot[title={refactorized every 150 updates}]
    \addplot[cred] table[x=it, y=stored, col sep=comma] {data/fill_I3_tau150.csv};
    \addplot[cblue] table[x=it, y=effective, col sep=comma] {data/fill_I3_tau150.csv};
    \addplot[cgrey, dashed] table[x=it, y=dense, col sep=comma] {data/fill_I3_tau150.csv};
\end{groupplot}
\end{tikzpicture}
\caption{Contents of the factor $L$ of $\Heff$ over a run (I3, equality variant), sampled
every 25 iterations: stored off-diagonal entries, the nonzero ones among them, and the
off-diagonal entries of a dense factor of the same dimension.}
\label{fig:fill}
\end{figure}

\paragraph{The true factor is dense.} In both panels the effective entries coincide with the
dense equivalent: the factor of $\Heff$ has essentially no exploitable sparsity. This
confirms \cref{rem:natural-ordering}: a fill-reducing ordering cannot help, and the cost of
a solve grows quadratically with $|\Nx|$.

\paragraph{Structural zeros accumulate.} Without the periodic refactorization, the stored
pattern only grows. On I3 it reaches almost a million entries in the column phase, more than
fifteen times the size of the actual factor, and almost all of them are structural zeros.
Every solve and every update pays for them, and they connect active indices with inactive
ones, which enlarges the reach of the sparse solves. The drop around iteration 1\,400 is a
refactorization triggered by a degraded pivot (\cref{sec:detection}). It resets the
pattern, after which it grows again and stays at about twice the size of the factor. With
refactorization every 150 updates, the stored pattern stays within a small margin of the
effective one throughout.

\paragraph{The refactorization interval.} \Cref{tab:refactor} and \cref{fig:refactor} vary the
refactorization interval. Without refactorization, the method is three to five times slower
and CHOLMOD needs six to eleven times more memory. Refactorizing too often is also expensive,
since every rebuild costs $O(|\Nx|^3)$ (\cref{prop:refactor-cost}): with an interval of 25, I3
spends most of its time on almost 3\,000 rebuilds. Between these extremes the optimum is flat.
Any interval between 100 and 600 is within about 10\% of the best time, and the default of 150
is within 4\% of the best on both instances.

\begin{table}[ht]
\centering
\small
\renewcommand{\arraystretch}{1.15}
\resizebox{\textwidth}{!}{%
\begin{tabular}{llrrrrrrr}
\toprule
 & interval $\tau$ & time [s] & refact. & max.\ stored & max.\ struct.\ zeros & mean zero share & reach$/|\Nx|$ & CHOLMOD peak [MB] \\
\midrule
\tabinput{data/tab_refactor.tex}
\bottomrule
\end{tabular}
}
\caption{Effect of the refactorization interval (equality variants). ``reach$/|\Nx|$'' is the
average number of rows touched by a sparse solve per nullspace column. ``never'' still
refactorizes when a pivot degrades.}
\label{tab:refactor}
\end{table}

\begin{figure}[ht]
\centering
\begin{tikzpicture}
\begin{axis}[width=0.68\textwidth, height=0.42\textwidth, xmode=log, xlabel={refactorization interval $\tau$},
    ylabel={total time [s]}, xtick={25,50,100,150,300,600,2000},
    xticklabels={25,50,100,150,300,600,never}, tick label style={font=\scriptsize},
    label style={font=\small}, legend style={font=\scriptsize, at={(0.02,0.98)}, anchor=north west},
    ymin=0, unbounded coords=jump, every axis plot/.append style={line width=0.8pt}]
    \addplot[cblue, mark=*, mark size=1.5pt] table[x=tau, y=time, col sep=comma] {data/refactor_I2.csv};
    \addplot[corange, mark=square*, mark size=1.5pt] table[x=tau, y=time, col sep=comma] {data/refactor_I3.csv};
    \legend{I2, I3}
\end{axis}
\end{tikzpicture}
\caption{Total time of the active-set method against the refactorization interval.}
\label{fig:refactor}
\end{figure}


\section{Discussion}
\label{sec:discussion}

The experiments support the main claims of this work.
\begin{itemize}
    \item \textbf{Correctness.} The active-set method solves the same sequence of master
          problems as a general-purpose QP solver, reproduces the optimum of the full linear
          program, and certifies it by the Lagrangian bound, for equality and inequality rows.
    \item \textbf{Speed.} A master problem of the active-set method costs between $0.01$ and
          $7$\,ms on average on our instances. That is two to three orders of magnitude less than a
          HiGHS QP master and two orders of magnitude less than a warm-started LP master. This
          makes regularized column generation much faster than classical column generation,
          although it often needs more iterations: on I3, 7.7\,s against 650\,s.
    \item \textbf{Inequality rows} cost nothing extra: the surplus operations are rank-one
          updates of the same factorization, the factorized matrix only shrinks when
          surpluses become free, and the runs are shorter.
    \item \textbf{Refactorization} is essential for long runs, and a fixed interval of a
          few hundred updates works well.
\end{itemize}

The experiments also show the limits of the approach. The factor of $\Heff$ is dense, so the
cost of the method grows with the cube of the number of linking rows, and the
\textbf{proximal phase} can be long, since the duals approach the restricted optimum in many
small steps (\cref{prop:center-only}). Two natural improvements are therefore a dense
factorization with cheap rank-one updates once $|\Nx|$ is large, and an adaptive
stabilization parameter in the proximal phase. A change of $\delta$ does not affect the
factorization (\cref{rem:factor-independent}) and is therefore free for this method.
